% Generated by tex/build_swebok_structure.py from tex/chapters/ch01/06_要求仕様化.tex
\subsubsection{受け入れ基準に基づく要求仕様化}

\term{受け入れ基準}とは、要求が満たされたと判断するための条件である。要求を受け入れ基準の形で書くと、テストと要求が強く結びつく。

\term{受け入れテスト駆動開発}（ATDD）は、機能を作る前に、その機能が完成したと判断するテストを合意する方法である。\term{振る舞い駆動開発}（BDD）は、利用者の振る舞いを「前提、事象、結果」の形で書く方法である。

\MiniBox{BDDの例}{前提：口座残高が500円で、カードが有効である。事象：利用者が100円の出金を要求する。結果：ATMは100円を払い出し、口座残高は400円になる。}

この方法は、曖昧さの低減に有効である。一方で、どのシナリオを書くべきかを見落とすと不完全性は残る。そのため、境界値分析、組合せテスト、同値分割などのテスト技法と組み合わせるとよい。
